\section{迁徙不是插曲，而是国家能力的尺度}
\label{sec:synthesis-migration}

从永嘉到侯景之乱，再到隋灭陈，迁徙始终改变着谁能被征发、谁可以任官、何处有军粮、哪一座城能成为都城。南渡的家庭不是一支整齐移动的士族队伍；侨置州郡也不是人口普查的同义词。官员、军人、工匠、僧侣、依附者和地方居民在不同时间抵达、停留或再迁徙，带来旧籍、职业、婚姻和土地要求。建康朝廷可以利用这些联系，又必须同江东既有社会协商。墓志的籍贯、官职和婚配能让我们看见少数有文字的家庭，却不能把他们的行迹变成全体流民的经验。

北方的移动也不该只用“胡汉”两字概括。前赵、后赵、前秦、诸燕、北魏和河西政权依靠不同的城郭、部众、旧郡县吏员、牧地和绿洲道路组织资源。史书的“降”“叛”“部”“胡”“羌”等词，首先是军政与叙事分类；它们在不同政权、不同战争和不同书写者那里未必指同一群人。最可靠的提问不是给每个称谓配一个不变的现代标签，而是追问它何时进入征发、安置、赏罚或边防的语境。

\section{宗教是公共资源网络，也是材料边界}
\label{sec:synthesis-religion}

寺院、道观、石窟、译场和造像把僧侣、供养者、工匠、木石、道路与城市市场连接起来。后秦长安的译经、北魏云冈和洛阳的营造、南朝建康的寺院网络、北齐河北的造像和北周的禁断，都不能只写成统治者的个人信仰。它们涉及谁能组织劳力、保存经卷、占有土地、承接施舍并在城市中获得可见性。相反，太武帝灭佛和周武帝禁佛也不只是思想立场：诏令把僧尼、寺产和器物纳入国家清查，却不能证明每一处寺院同时被毁，或每一个人的信仰在一夜间改变。

宗教材料尤其要求区分层次。僧传和护教文本擅长保存典范人物与神异意义；经录可提示文本流传；造像题记能校验特定施主、日期和地点；遗址显示工程曾经发生。没有一类材料能自动说明所有居民如何信仰，或寺院在某州究竟有多少田产。把不同材料的沉默保留下来，并不削弱宗教史，反而让寺院和道团得以作为战争、城市和行政中真实而有限的行动者出现。

\section{边境不是边缘：河西、草原、淮汉与海路}
\label{sec:synthesis-frontier}

西晋的统一依靠长江水陆走廊；东晋的存续取决于江淮、荆州和海路能否输送人粮；北魏面对柔然、河西与高昌时，则把军镇、牧地、使节和绿洲城市放进同一套压力之中。五、六世纪北齐、北周和隋与突厥的关系进一步说明，北方政权不能只用中原的税粮决定战争节奏。婚姻、贡赐、降附、互市和边镇供给，都是政治选择的条件。

因此，本卷的示意图不画精确疆界。河流、山道、城郭和仓储往往比色块边缘更能解释一支军队为何能到达、为何不能久守。战争史料中的兵数、斩获和路线常有夸饰；城市、道路、文书与地方材料可以约束想象，却也不能独自复原一场战役。边境是人、货、文本和消息来回穿过的区域，不是给中原帝国提供舞台布景的静止外圈。

\section{行政：从诏令到地方，永远隔着中介}
\label{sec:synthesis-administration}

西晋律令、北魏的班禄均田三长、东西魏北周的军府、隋的州郡户籍与租调，都是使人口、土地、劳役和军队可被登记、比较和调动的尝试。\eventanchor{NSL-440-589-073}{北魏·均田令}与\eventanchor{NSL-440-589-075}{北魏·三长制}可以证明四八五至四八六年制度被颁布，也能显示朝廷试图把土地分类、邻里责任、赋役和治安接合起来；它们不能证明丈量、授受和征税在所有地区同步发生。宗族、旧部落、坞壁、寺院、州郡吏员和地方强宗，既可能成为制度中介，也可能使一纸命令改变形状。

“汉化”“鲜卑化”同样不能替代行政史。\eventanchor{NSL-440-589-086}{北魏·迁都洛阳}四九四年北魏迁都洛阳后，改姓、服饰、礼仪和婚姻政策显示洛阳朝廷倡导的公开规范，却不足以测量各地语言、家庭关系或身份认同。北魏裂解后，东魏北齐的河北网络与西魏北周的关中网络各自吸纳军人、流民和地方精英；这正是制度从来不能脱离地理与中介的结果。到五八九年，隋接收陈的州郡、仓储、官员、宗族和寺院，仍须面对同一个问题：最高权力的转换能否化为地方可持续的登记、供给和裁决。

\begin{sourcenote}[综合章的材料责任]
本章以正史志书和编年材料确定政策、战争与政权更替的最低骨架，以墓志、造像、佛道文本、河西与吐鲁番文书、城址和道路遗存检查它们在何处可能触及地方。后者并非“更真实”的社会全景，前者也并非天然不可信；它们回答的问题不同。材料相遇处可以提出解释，不能相遇处则必须保留为空白。
\end{sourcenote}
